
\section{Automatic trajectories and a critical-owner exchange hypothesis}
\subsection{Automatic outcomes and the proposal-class limit}
All four registered actors finish under their original versions and rules. Both Tf policies accept the selected thirteen-connection result, reaching 13.413972155262854 ns with wire 2074301, SB105 and CB5062. Each then selects from that accepted current state and finds no positive group in the current operator class. Exact identical evidence is reused explicitly; different historical creation costs are not reassigned as new algorithm runtime.

A0 strict acceptance performs two automatically selected transactions: CPD changes from 13.607436066820355 to 13.421610489672275 and then to 13.413972155262854 ns. The second group is selected from the first accepted state's fresh feedback, not from a fixed initial list. Final wire is 2070933, SB98 and CB5067. The earlier all-positive single transaction reaches the same CPD with twelve fewer wire units. Thus automatic target updating is operational, but this two-step schedule does not improve the four-axis quality frontier. Complete comparable algorithm-time dominance is unknown; imported-state actor times exclude referenced historical creation, initial routing and static caches.

The A0 endpoint actor rejects its first selected group and retains the pristine state. This is a limitation of the positive-winning-path proposal class, rather than evidence that endpoint acceptance has no opportunity: two already legal owner-only alternatives from the identical current state pass exact endpoint-vector acceptance. The weighted alternative also reduces wire by eight; the physical alternative increases wire by thirty-two. A replay check recomputes acceptance against the same initial budgets and current endpoint/domain identities without new native work. These references motivate a future proposal portfolio; they are not a history-dependent algorithm or per-circuit best-result selection. A rejected primary proposal must not be interpreted as exhausting resource/vector proposals.

\subsection{Conditional critical-owner competition}
Fresh feedback from the accepted Tf state still finds a 185.825799193 ps local relaxed gain for target 26319/pin1, blocked in part by critical owner 26320. Automatically releasing the original target-witness conflicts identifies seven connections. After fixing that target witness, independent exact capacity queries require the critical owner's pin1 to increase by 549.838397035 ps. Both current winning paths then have an optimistic net increase of 364.012597842 ps. This rules out a gain for that particular constrained local model, not every target path or larger transaction.

Relaxing other noncritical background occupancy permits owner paths with their original physical delays. Following those witnesses identifies six additional noncritical branches, expanding the proposal to thirteen connections. The core target and critical-owner paths can coexist under the resulting partial occupancy. Missing noncritical pins still require restoration; this is not a complete legal routing or an STA result. The kernel remains native search. The new hypothesis concerns dependency closure, target-witness adoption and scheduling, rather than a new primitive whose benefit has already been demonstrated.

A separately registered five-arm native diagnosis contrasts minimal/expanded closure with a fixed target witness, expanded native free target choice, physical-delay/hard-capacity treatment and owner-first scheduling. Native free control has the same source, full-device base domain, branch scope and one acceptance opportunity, with a wider target-path choice than the fixed-witness arm. The objectives change cost/capacity treatment together. Every arm must restore original state, preserve retained branches and terminals, satisfy explicit preset edges, and pass complete physical legality/native FULL/STA plus exact initial-budget decisions. These current-repaired-state diagnostics do not replace future automatic evaluation from pristine strong starts.

\subsection{Exact conditional distance reuse}
No queries are repeated during scope expansion. The previously exact noncritical-relaxed result supplies a lower bound because its domain includes the expanded actual-capacity query. Its same path remains feasible after the extra owner branches are released, supplying an equal upper bound. Retained nodes, anchor delays, graph/core semantics and explicit blocked-set inclusion are checked before issuing this conditional equality certificate. Independent shortest-path comparisons on 96 small random DAGs produce 95 valid certificates and one blocked-witness miss; changed anchors, targets, unfinished bounds and invalid domain inclusions are rejected in 371 controls. These checks establish the tested path-certificate behavior, not a global FPGA optimum, native floating-point STA bound or final acceptance. Scope expansion uses zero new solver queries or route commits.
